% GENERATED FILE. Edit canonical lessons, then run npm run generate:books.
% canonical-source-hash: fnv1a64:7e7eded69c4d1746
% canonical-lessons: IT-C01-ciao, IT-W01-ciao-guided-copy, IT-C01-buono, IT-C01-il-la-lo, IT-C01-giorno, IT-C01-buongiorno, IT-C01-sera, IT-C01-notte, IT-C01-grazie, IT-C01-practice

\chapter{Greetings}
\label{ch:greetings}
\hlchaptermodality{1}

\begin{chapteropening}
\textbf{By the end of this chapter:} \emph{I can greet someone in Italian, recognize the reply, and choose an appropriate response.}

Complete IT-C01-practice, the canonical closing checkpoint, and demonstrate the chapter promise: greet someone in Italian, recognize the reply, and choose an appropriate response.
\end{chapteropening}

\section[ciao]{ciao — \textquotedblleft{}hi,\textquotedblright{} and the word that means \textquotedblleft{}I am your slave\textquotedblright{}}

\label{lesson:IT-C01-ciao}



\begin{quote}
The most famous \textquotedblleft{}hello\textquotedblright{} in the world — and it hides an astonishing history.
\end{quote}

\begin{sounds}
\begin{itemize}
\raggedright
  \item \textbf{ciao} = \emph{CHOW}: \textbf{ci} before a vowel is the \emph{ch} of \emph{church}; \emph{-ao} is one gliding syllable.
\end{itemize}
\end{sounds}

\begin{cousinweb}
\textbf{ciao} began as the Venetian greeting \emph{s-ciào (vostro)} — \textquotedblleft{}{[}I am your{]} slave / servant,\textquotedblright{} a humble \textquotedblleft{}at your service.\textquotedblright{} That \emph{s-ciào} is Venetian for \textbf{schiavo} (\textquotedblleft{}slave\textquotedblright{}) $\leftarrow$ Medieval Latin \textbf{sclavus}. The very same word is English \textbf{slave} — and, astonishingly, \textbf{Slav}: medieval slaves were so often Slavic captives that the ethnic name \emph{became} the word for \textquotedblleft{}slave\textquotedblright{} across Europe. So every casual \emph{ciao} is a fossil of \textquotedblleft{}I am your servant.\textquotedblright{}
\end{cousinweb}

\begin{culture}
\emph{ciao} is \textbf{informal}, and works both coming and going (\textquotedblleft{}hi\textquotedblright{} and \textquotedblleft{}bye\textquotedblright{}) — for friends, family, peers. To someone you'd show respect, or on a first meeting, you reach for a polite daytime greeting instead. You will build that greeting from two small pieces later in this chapter. Doubled, \emph{ciao ciao}, it's a breezy \textquotedblleft{}bye-bye.\textquotedblright{}
\end{culture}

\subsection*{Writing — follow the greeting you can see}
Keep \textbf{ciao} visible. Point to \textbf{ci}, the two letters that begin the \emph{ch} sound, then slide your finger through \textbf{ao}. Trace the whole word once from \textbf{c} to \textbf{o}, and say \emph{CHOW}.

This is not a spelling test. Your eye and hand are only following one short greeting whose sound and meaning are still on the page.

\subsection*{Your turn}
Say these aloud:

\begin{itemize}
\raggedright
  \item \textquotedblleft{}ciao\textquotedblright{} (\emph{CHOW})
  \item its buried meaning — \textquotedblleft{}I am your slave/servant\textquotedblright{}
  \item the English cousins — \emph{slave}, \emph{Slav}
\end{itemize}

\begin{tcolorbox}[breakable,colback=teal!4,colframe=teal!35!black,title={Before you move on}]
What did \textbf{ciao} literally mean, and what two English words share its root? (\textquotedblleft{}I am your slave\textquotedblright{}; \emph{slave} and \emph{Slav}, from Latin \emph{sclavus}.)
\end{tcolorbox}
\section[ciao]{ciao — copy one greeting with help}

\label{lesson:IT-W01-ciao-guided-copy}



\begin{quote}
Point to \textbf{ciao} and say it once: \emph{CHOW}. The model stays on the page.
\end{quote}

\subsection*{Writing — one guided copy}
Copy \textbf{ciao} once on the line below it. Write \textbf{ci}, pause, then add \textbf{ao}. Look back at the model whenever you want. Check that the \textbf{i} has its dot.

Say the word as you finish. The four letters record the sound you already know; there is nothing to guess and no second word today.

\begin{tcolorbox}[breakable,colback=teal!4,colframe=teal!35!black,title={Before you move on}]
Compare your copy with the model. If one letter differs, repair only that letter. No hiding, no spelling from memory, and no free sentence yet.
\end{tcolorbox}
\section[buono / buon]{buono / buon — \textquotedblleft{}good\textquotedblright{}}

\label{lesson:IT-C01-buono}



\begin{quote}
One small word meaning \textquotedblleft{}good.\textquotedblright{} Learn it alone first; later lessons will show how it helps build several greetings.
\end{quote}

\begin{sounds}
\begin{itemize}
\raggedright
  \item \textbf{buono} = \emph{BWOH-no} (the \emph{uo} is a glide). Before a masculine noun it clips to \textbf{buon} (\emph{bwon}).
\end{itemize}
\end{sounds}

\begin{cousinweb}
\textbf{buono} (\textquotedblleft{}good\textquotedblright{}) $\leftarrow$ Latin \textbf{bonus} $\to$ English \textbf{bon}us, \textbf{bon}anza, \textbf{bon}bon, and (via French) \textbf{boun}ty. Spanish, another daughter of Latin, made \emph{bonus} into \emph{bueno}; Italian kept it closer, \emph{buono}.
\end{cousinweb}

\begin{grammarlens}[title={Grammar Lens: it agrees with the noun}]
\textbf{buono} changes to match its noun's gender and number: \emph{buono} (m.), \emph{buona} (f.), \emph{buoni} / \emph{buone} (pl.) — Latin's legacy across the Romance family. Before a masculine noun it can drop the \emph{-o} to \textbf{buon}. For now, just hear and say that small change; the next nouns will give it a job.
\end{grammarlens}

\subsection*{Your turn}
Say these aloud:

\begin{itemize}
\raggedright
  \item \textquotedblleft{}buono\textquotedblright{}, then clipped \textquotedblleft{}buon\textquotedblright{}
  \item English cousins — \emph{bonus}, \emph{bonanza}
\end{itemize}

\begin{tcolorbox}[breakable,colback=teal!4,colframe=teal!35!black,title={Before you move on}]
What Latin word is \textbf{buono} from, and what short form can it take before a masculine noun? (\emph{bonus}; \emph{buon}, with the final \emph{-o} dropped.)
\end{tcolorbox}
\section[il / la / lo]{il / la / lo — \textquotedblleft{}the,\textquotedblright{} and the gender it carries}

\label{lesson:IT-C01-il-la-lo}



\begin{quote}
Before your first noun: the words for \textquotedblleft{}the.\textquotedblright{} Italian splits \textquotedblleft{}the\textquotedblright{} by gender — a habit inherited from Latin.
\end{quote}

\begin{cousinweb}
\textbf{il} and \textbf{lo} (masculine \textquotedblleft{}the\textquotedblright{}) and \textbf{la} (feminine) $\leftarrow$ Latin \textbf{ille} / \textbf{illa} (\textquotedblleft{}that one\textquotedblright{}), the pointing word worn down to \textquotedblleft{}the.\textquotedblright{} The same \emph{ille/illa} gave Spanish \emph{el/la} and French \emph{le/la}. English \textquotedblleft{}the\textquotedblright{} is Germanic and unrelated.
\end{cousinweb}

\begin{grammarlens}[title={Grammar Lens: gender, and two masculine \textquotedblleft{}the\textquotedblright{}s}]
Every Italian noun is \textbf{masculine or feminine} — Latin's three genders (the neuter folding into the masculine) collapsed to two. Learn each noun \emph{with} its article. Italian has \textbf{two} masculine articles: \textbf{lo} before tricky clusters at the start of a noun, \textbf{il} in the ordinary pattern; feminine is always \textbf{la}. Your first noun comes next, so you can practise one real pairing then.
\end{grammarlens}

\subsection*{Your turn}
Say these aloud:

\begin{itemize}
\raggedright
  \item \textquotedblleft{}il\textquotedblright{}, \textquotedblleft{}la\textquotedblright{}, \textquotedblleft{}lo\textquotedblright{}
  \item their source — Latin \emph{ille / illa}
\end{itemize}

\begin{tcolorbox}[breakable,colback=teal!4,colframe=teal!35!black,title={Before you move on}]
What Latin word gave \textbf{il / lo / la}, and how many genders does Italian have? (\emph{ille / illa}; two — masculine and feminine.)
\end{tcolorbox}
\section[giorno]{giorno — \textquotedblleft{}day,\textquotedblright{} and the road to English \emph{journey}}

\label{lesson:IT-C01-giorno}



\begin{quote}
Your first noun — masculine, and a cousin of \emph{journal} and \emph{journey}.
\end{quote}

\begin{sounds}
\begin{itemize}
\raggedright
  \item \textbf{gi} = the \emph{j} of \emph{jam}; \textbf{giorno} = \emph{JOR-no}.
\end{itemize}
\end{sounds}

\begin{cousinweb}
\textbf{il giorno} (masc.) $\leftarrow$ Latin \textbf{diurnum} (\textquotedblleft{}daytime\textquotedblright{}), from \textbf{dies} (\textquotedblleft{}day\textquotedblright{}). That detour through \emph{diurnum} is why \emph{giorno} looks nothing like \emph{dies}:

\begin{itemize}
\raggedright
  \item via \emph{diurnum}: \textbf{jour}nal, ad\textbf{journ}, so\textbf{journ}, \textbf{journ}ey (a \emph{day's} travel);
  \item via \emph{dies}: \textbf{di}urnal, \textbf{di}ary.
\end{itemize}

Italian \emph{giorno} and French \emph{jour} are near-twins here; Spanish, taking \emph{dies} directly, got the very different \emph{día}.
\end{cousinweb}

\begin{grammarlens}[title={Grammar Lens: masculine, plural by vowel-change}]
\emph{giorno} is \textbf{masculine} (\emph{il giorno}). Plural \textbf{i giorni} — Italian marks the plural by changing the final vowel (\emph{-o} $\to$ \emph{-i}), not by adding \emph{-s}.
\end{grammarlens}

\subsection*{Your turn}
Say these aloud:

\begin{itemize}
\raggedright
  \item \textquotedblleft{}il giorno\textquotedblright{} (masc.), \textquotedblleft{}i giorni\textquotedblright{} (pl.)
  \item the cousins — \emph{journal}, \emph{journey}, \emph{diary}
\end{itemize}

\begin{tcolorbox}[breakable,colback=teal!4,colframe=teal!35!black,title={Before you move on}]
What Latin word gave \textbf{giorno}, and what English words share it? (\emph{diurnum} $\leftarrow$ \emph{dies}; \emph{journal}, \emph{journey}, \emph{diurnal}, \emph{diary}.) How does Italian form the plural? (Changes the final vowel: \emph{-o} $\to$ \emph{-i}.)
\end{tcolorbox}
\section[buongiorno]{buongiorno — \textquotedblleft{}good day,\textquotedblright{} assembled}

\label{lesson:IT-C01-buongiorno}



\begin{quote}
Put the two pieces together into the polite, all-purpose hello.
\end{quote}

\begin{cousinweb}
\begin{itemize}
\raggedright
  \item buon (\textquotedblleft{}good\textquotedblright{}) + giorno (\textquotedblleft{}day\textquotedblright{}).
\end{itemize}

\begin{quote}
\textbf{buongiorno} = \textquotedblleft{}good day,\textquotedblright{} one word.
\end{quote}
\end{cousinweb}

\begin{culture}
Where the casual \emph{ciao} says \textquotedblleft{}I'm your servant,\textquotedblright{} \emph{buongiorno} simply wishes you a good day — and it's the \textbf{safe, respectful} choice with anyone: shopkeepers, strangers, elders. Use it through the day until late afternoon. A later lesson will build the corresponding evening greeting.
\end{culture}

\subsection*{Your turn}
Say these aloud:

\begin{itemize}
\raggedright
  \item \textquotedblleft{}buongiorno\textquotedblright{} — buon + giorno
  \item when to use it (respect, daytime) vs \emph{ciao} (friends)
\end{itemize}

\begin{tcolorbox}[breakable,colback=teal!4,colframe=teal!35!black,title={Before you move on}]
What two words make \textbf{buongiorno}, and when do you use it rather than \emph{ciao}? (\emph{buon} + \emph{giorno}; for respect / strangers / daytime.)
\end{tcolorbox}
\section[sera / buonasera]{sera / buonasera — \textquotedblleft{}evening\textquotedblright{} and \textquotedblleft{}good evening\textquotedblright{}}

\label{lesson:IT-C01-sera}



\begin{quote}
The evening swap for \emph{buongiorno}.
\end{quote}

\begin{cousinweb}
\textbf{la sera} (fem.) $\leftarrow$ Latin \textbf{sera} (\textquotedblleft{}late {[}hour{]}\textquotedblright{}), from \textbf{serus}, \textquotedblleft{}late.\textquotedblright{} English cousin: the rare \emph{serotine} (a bat of the \emph{late} evening). \textbf{buonasera} = \emph{buona} + \emph{sera}, \textquotedblleft{}good evening\textquotedblright{} — \emph{buona} in its \textbf{feminine} form, because \emph{sera} is feminine (contrast the masculine \emph{buon giorno}).
\end{cousinweb}

\begin{grammarlens}[title={Grammar Lens: feminine agreement}]
Notice \emph{buon\textbf{a} sera} vs. \emph{buon giorno}: the \textquotedblleft{}good\textquotedblright{} word matches its noun's gender — feminine \emph{sera} takes \emph{buona}, masculine \emph{giorno} takes \emph{buon}. Same adjective, two coats.
\end{grammarlens}

\subsection*{Your turn}
Say these aloud:

\begin{itemize}
\raggedright
  \item \textquotedblleft{}buonasera\textquotedblright{} (\emph{buona} + \emph{sera})
  \item why it's \emph{buona} here but \emph{buon} in \emph{buon giorno} (gender)
\end{itemize}

\begin{tcolorbox}[breakable,colback=teal!4,colframe=teal!35!black,title={Before you move on}]
What Latin word gave \textbf{sera}, and why is it \emph{buonasera} but \emph{buon giorno}? (\emph{sera} $\leftarrow$ \emph{serus}, \textquotedblleft{}late\textquotedblright{}; \emph{sera} is feminine so \textquotedblleft{}good\textquotedblright{} is \emph{buona}.)
\end{tcolorbox}
\section[notte / buonanotte]{notte / buonanotte — \textquotedblleft{}night,\textquotedblright{} and the \emph{-ct-} $\to$ \emph{-tt-} rule}

\label{lesson:IT-C01-notte}



\begin{quote}
\textquotedblleft{}Night\textquotedblright{} — and a sound-rule that decodes dozens of Italian words at once.
\end{quote}

\begin{cousinweb}
\textbf{la notte} (fem.) $\leftarrow$ Latin \textbf{noctem} (\emph{noct-}) $\to$ English \textbf{noct}urnal, equi\textbf{nox}. Watch just one change today: Latin \emph{-ct-} became Italian \textbf{-tt-}, so \emph{noc\textbf{t}em} became \emph{no\textbf{tt}e}. Spanish and French reshaped that same root in their own ways. Later words will let you see whether this pattern repeats; there is no list to memorize now.

\textbf{buonanotte} = \emph{buona} + \emph{notte}, \textquotedblleft{}good night\textquotedblright{} — the farewell for bed (not an evening hello; use the evening greeting you already learned for that).
\end{cousinweb}

\subsection*{Your turn}
Say these aloud:

\begin{itemize}
\raggedright
  \item \textquotedblleft{}buonanotte\textquotedblright{}
  \item the one change you saw — Latin \emph{noctem} $\to$ Italian \emph{notte}
\end{itemize}

\begin{tcolorbox}[breakable,colback=teal!4,colframe=teal!35!black,title={Before you move on}]
What Latin word gave \textbf{notte}, and what regular sound-change does it show? (\emph{noctem}; Latin \emph{-ct-} $\to$ Italian \emph{-tt-}.) When is \emph{buonanotte} used? (As a bedtime farewell.)
\end{tcolorbox}
\section[grazie]{grazie — \textquotedblleft{}thank you,\textquotedblright{} from \emph{grace}}

\label{lesson:IT-C01-grazie}



\begin{quote}
\textquotedblleft{}Thanks\textquotedblright{} — and it literally offers someone \emph{grace}.
\end{quote}

\begin{sounds}
\begin{itemize}
\raggedright
  \item \textbf{grazie} = \emph{GRAH-tsyeh}: \textbf{z} in Italian is \emph{ts}; the final \emph{-e} is sounded (so \emph{-zie} is two beats).
\end{itemize}
\end{sounds}

\begin{cousinweb}
\textbf{grazie} (\textquotedblleft{}thanks\textquotedblright{}) $\leftarrow$ Latin \textbf{gratiae}, plural of \textbf{gratia} (\textquotedblleft{}favour, thanks, grace\textquotedblright{}). It's the \emph{same} \emph{gratia} behind English \textbf{grace}, \textbf{grat}itude, \textbf{grat}eful, \textbf{grat}is (\textquotedblleft{}free\textquotedblright{}), con\textbf{grat}ulate, and in\textbf{grat}iate. So \emph{grazie} literally wishes someone \textbf{grace}.
\end{cousinweb}

\begin{culture}
Think of \textbf{grazie} as offering someone \textquotedblleft{}graces\textquotedblright{} for what they did. In the next chapter, you will learn the short reply people use after hearing it.
\end{culture}

\subsection*{Your turn}
Say these aloud:

\begin{itemize}
\raggedright
  \item \textquotedblleft{}grazie\textquotedblright{} (\emph{GRAH-tsyeh})
  \item the cousins — \emph{grace}, \emph{gratitude}, \emph{gratis}
\end{itemize}

\begin{tcolorbox}[breakable,colback=teal!4,colframe=teal!35!black,title={Before you move on}]
What Latin word is \textbf{grazie} from, and three English cousins? (\emph{gratia}, \textquotedblleft{}grace\textquotedblright{}; \emph{grace}, \emph{gratitude}, \emph{gratis}.)
\end{tcolorbox}
\section[Practice]{Chapter 1 — the day of greetings}

\label{lesson:IT-C01-practice}



\begin{quote}
No new words today. Everything below is something you have already met, said once, and can now put in one place.
\end{quote}

\begin{grammarlens}[title={What you've built — the day of greetings}]
\noindent
\begin{tabularx}{\linewidth}{@{}>{\raggedright\arraybackslash}X>{\raggedright\arraybackslash}X@{}}
\toprule
\textbf{Italian} & \textbf{English} \\
\midrule
\emph{ciao} & hi / bye (informal) \\
\emph{buongiorno} & good day \\
\emph{buonasera} & good evening \\
\emph{buonanotte} & good night \\
\emph{grazie} & thank you \\
\bottomrule
\end{tabularx}

Every piece from a Latin root:

\begin{itemize}
\raggedright
  \item \textbf{ciao} $\leftarrow$ \emph{sclavus} (\textquotedblleft{}slave\textquotedblright{}; English \emph{slave}, \emph{Slav})
  \item \textbf{buono} $\leftarrow$ \emph{bonus}
  \item \textbf{giorno} $\leftarrow$ \emph{diurnum} (English \emph{journal}, \emph{journey})
  \item \textbf{notte} $\leftarrow$ \emph{noctem}, via the \emph{-ct-} $\to$ \emph{-tt-} rule (\emph{notte}, not \emph{noche}/\emph{nuit})
  \item \textbf{grazie} $\leftarrow$ \emph{gratia} (\textquotedblleft{}grace\textquotedblright{}; English \emph{grace}, \emph{gratitude})
\end{itemize}

And two Italian habits: adjectives \textbf{agree} with the noun's gender (\emph{buon giorno} m. / \emph{buona sera} f.), and the plural is made by \textbf{changing the final vowel} (\emph{giorno} $\to$ \emph{giorni}), not adding \emph{-s}.
\end{grammarlens}

\subsection*{Your turn}
Say these aloud:

\begin{itemize}
\raggedright
  \item all five greetings
  \item greet and thank — \textquotedblleft{}buongiorno!\textquotedblright{} … \textquotedblleft{}grazie.\textquotedblright{}
  \item the buried story of \emph{ciao} (\textquotedblleft{}I am your slave\textquotedblright{})
\end{itemize}

\begin{tcolorbox}[breakable,colback=teal!4,colframe=teal!35!black,title={Before you move on}]
Say all five greetings. Which one once meant \textquotedblleft{}I am your slave\textquotedblright{}? (\emph{ciao}.) What sound-rule turns Latin \emph{noctem} into Italian \emph{notte}? (\emph{-ct-} $\to$ \emph{-tt-}.)

Next chapter, you will learn how to introduce yourself and how to address someone in a familiar or respectful way. Those new Italian forms can wait for their own small lessons.
\end{tcolorbox}
